% =====================================================================
%  TA-IND-01 - Ingenieria de Requerimientos (ISR-401)
%  Facultad de Ciencias de la Computacion - UTEQ
%  Estudiante: Jean Pierre Robinson
%
%  ARCHIVO PRINCIPAL. Compilar con pdfLaTeX + Biber:
%      pdflatex main.tex
%      biber main
%      pdflatex main.tex
%      pdflatex main.tex
%
%  Antes de entregar:
%   1) Reemplazar TODAS las marcas [PENDIENTE: ...].
%   2) Borrar los comentarios que empiezan con "GUIA".
%   3) Colocar las 3 capturas en figuras/ con los nombres indicados.
%   4) Actualizar la Seccion de declaracion de uso de IA.
% =====================================================================
\documentclass[12pt,a4paper]{article}

\usepackage[T1]{fontenc}
\usepackage[utf8]{inputenc}
\usepackage{lmodern}
\usepackage[spanish,es-noshorthands]{babel}
\usepackage{geometry}
\geometry{margin=2.5cm}
\usepackage{microtype}
\usepackage{graphicx}
\usepackage{booktabs}
\usepackage{longtable}
\usepackage{array}
\usepackage{tabularx}
\usepackage{enumitem}
\usepackage{xcolor}
\usepackage{csquotes}
\usepackage{caption}
\usepackage[hidelinks]{hyperref}
\usepackage[backend=biber,style=ieee]{biblatex}
\addbibresource{referencias.bib}

% ---------------------------------------------------------------------
% Ayudas de trabajo (no forman parte del contenido entregable)
% ---------------------------------------------------------------------
\newcommand{\pendiente}[1][]{\textcolor{red}{\textbf{[PENDIENTE#1]}}}

% Marcador compacto para celdas estrechas de tabla.
\newcommand{\pd}{\textcolor{red}{\scriptsize\textbf{PEND}}}

% Figura que compila aunque el archivo de la captura no exista todavia.
% Uso: \captura{archivo.png}{pie de figura}{etiqueta}
\newcommand{\captura}[3]{%
  \begin{figure}[htbp]
    \centering
    \IfFileExists{figuras/#1}{%
      \includegraphics[width=0.82\linewidth]{figuras/#1}%
    }{%
      \fbox{\parbox[c][5.5cm][c]{0.82\linewidth}{%
        \centering \pendiente{: reemplazar por la captura propia
        \texttt{figuras/#1}}}}%
    }%
    \caption{#2}
    \label{#3}
  \end{figure}%
}

\captionsetup{font=small,labelfont=bf}
\setlength{\parindent}{0pt}
\setlength{\parskip}{0.45em}
\renewcommand{\arraystretch}{1.25}

% ---------------------------------------------------------------------
\begin{document}

\begin{center}
  {\Large\bfseries Requisitos de un sistema conocido a partir de\\[2pt]
  ISO/IEC/IEEE 29148:2018}\\[10pt]
  {\large TA-IND-01 - Trabajo autónomo individual}\\[6pt]
  Ingeniería de Requerimientos (ISR-401) - Cuarto nivel - Paralelo A\\
  Facultad de Ciencias de la Computación - Universidad Técnica Estatal de Quevedo\\[6pt]
  Jean Pierre Robinson\\
  7 de noviembre de 2026
\end{center}

\vspace{0.5em}
\hrule
\vspace{0.5em}

% =====================================================================
\section{Descripción del sistema}
\label{sec:sistema}

% GUIA (borrar al terminar). Un solo párrafo de 150 a 200 palabras:
%   * Qué problema resuelve la aplicación y para qué existe.
%   * Quiénes son los usuarios y qué hace cada rol (pasajero, conductor,
%     administrador u otros que realmente uses).
%   * En qué contexto se usa: dispositivo, conectividad, movilidad,
%     datos personales y de ubicación involucrados.
%   * No describas pantallas: de eso se encargan las figuras.
% Recuerda: el sistema debe estar registrado en el foro del SGA.

\pendiente{: párrafo de descripción del sistema (propósito, usuarios y
contexto de uso).}

Las Figuras~\ref{fig:cap1}, \ref{fig:cap2} y \ref{fig:cap3} son capturas
propias de la aplicación y respaldan los requisitos de la
Sección~\ref{sec:requisitos}.

\captura{captura-01.png}{\pendiente{: pie de la Figura 1; indicar qué se
observa y qué requisito respalda.}}{fig:cap1}

\captura{captura-02.png}{\pendiente{: pie de la Figura 2; indicar qué se
observa y qué requisito respalda.}}{fig:cap2}

\captura{captura-03.png}{\pendiente{: pie de la Figura 3; indicar qué se
observa y qué requisito respalda.}}{fig:cap3}

% =====================================================================
\section{Resumen analítico de ISO/IEC/IEEE 29148:2018}
\label{sec:resumen}

% GUIA (borrar al terminar).
%   * Extensión conjunta de 2.1 a 2.5: entre 1,5 y 2 páginas.
%   * Es un resumen ANALITICO, no una traducción ni una lista copiada:
%     explica con tus palabras, relaciona las ideas entre sí y muestra
%     por qué importan al escribir requisitos.
%   * Cada párrafo debe cerrar con la cláusula que lo respalda,
%     por ejemplo: "(5.2.5)".
%   * Prohibido pegar texto traducido de la norma sin interpretación.

Fuentes de este apartado: \cite{iso29148} (cláusulas 1, 3 y 5.2), con
apoyo en \cite{pohl2025} y \cite{swebok2024}.

\subsection{Qué es un requisito bien formado (5.2.4)}
\label{sec:bien-formado}

% GUIA: explica la CONSTRUCCION que propone 5.2.4, no solo la definición.
% Preguntas para orientar el párrafo:
%   * Qué elementos debe contener el enunciado (condición, sujeto,
%     acción, objeto, restricción) y para qué sirve cada uno.
%   * Por qué un requisito construido así es verificable.
%   * Da un ejemplo breve, mal formado y luego bien formado.

\pendiente{: explicación de la construcción del requisito bien formado
según 5.2.4, con un ejemplo propio.}

\subsection{Características de un requisito individual (5.2.5)}
\label{sec:individual}

% GUIA: las NUEVE características. Para cada una: qué significa y cómo se
% comprueba en la práctica. No basta enumerarlas; explica el criterio de
% verificación. Cierra cada idea con (5.2.5).
% Verifica que no te falte ninguna antes de entregar.

\pendiente{: explicación de las nueve características de un requisito
individual y de cómo se comprueba cada una.}

\subsection{Características de un conjunto de requisitos (5.2.6)}
\label{sec:conjunto}

% GUIA: qué debe cumplir el CONJUNTO (no el requisito aislado). Explica
% cada característica y pon un ejemplo de conflicto entre dos requisitos
% del sistema elegido.

\pendiente{: explicación de las características del conjunto de
requisitos según 5.2.6.}

\subsection{Criterios de lenguaje y términos prohibidos (5.2.7)}
\label{sec:lenguaje}

% GUIA: explica los criterios de redacción y presenta una lista de al
% menos diez términos ambiguos que la norma desaconseja, con la razón por
% la que cada uno impide verificar el requisito.

\pendiente{: criterios de lenguaje de 5.2.7 y lista razonada de términos
prohibidos.}

\subsection{Atributos de los requisitos (5.2.8)}
\label{sec:atributos}

% GUIA: distingue con claridad ATRIBUTO de TIPO de requisito. Un atributo
% no describe lo que el sistema hace, sino cómo se gestiona el requisito
% (identificador, prioridad, estado, fuente, criterio de verificación).

\pendiente{: explicación de los atributos de 5.2.8 y distinción entre
atributo y tipo de requisito.}

% =====================================================================
\section{Glosario de términos}
\label{sec:glosario}

% GUIA: diez términos tomados de la clausula 3 de la norma. La definición
% debe estar PARAFRASEADA (no copiada) y el ejemplo debe ser propio, del
% sistema elegido. Si un término no aparece en la clausula 3, no sirve.

\begin{longtable}{@{}p{3.0cm}p{6.3cm}p{5.0cm}@{}}
\caption{Términos de la cláusula 3 aplicados al sistema elegido.}
\label{tab:glosario}\\
\toprule
\textbf{Término} & \textbf{Definición parafraseada} &
\textbf{Ejemplo en el sistema} \\
\midrule
\endfirsthead
\multicolumn{3}{@{}l}{\footnotesize\itshape Continuación de la Tabla~\ref{tab:glosario}}\\
\toprule
\textbf{Término} & \textbf{Definición parafraseada} &
\textbf{Ejemplo en el sistema} \\
\midrule
\endhead
\bottomrule
\endfoot
T01 & \pendiente{} & \pendiente{} \\
T02 & \pendiente{} & \pendiente{} \\
T03 & \pendiente{} & \pendiente{} \\
T04 & \pendiente{} & \pendiente{} \\
T05 & \pendiente{} & \pendiente{} \\
T06 & \pendiente{} & \pendiente{} \\
T07 & \pendiente{} & \pendiente{} \\
T08 & \pendiente{} & \pendiente{} \\
T09 & \pendiente{} & \pendiente{} \\
T10 & \pendiente{} & \pendiente{} \\
\end{longtable}

% =====================================================================
\section{Requisitos del sistema}
\label{sec:requisitos}

% GUIA (borrar al terminar).
%   * Quince requisitos (R01 a R15) del sistema de la Seccion 1.
%   * Al menos CINCO tipos distintos de la clausula 5.2 (por ejemplo:
%     funcional, de desempeno, de interfaz externa, de proceso, de
%     calidad, de usabilidad o calidad en uso, de factores humanos,
%     restriccion). Verifica que cada tipo este realmente representado.
%   * Al menos CINCO requisitos de calidad, clasificados con
%     caracteristica Y subcaracteristica de ISO/IEC 25010:2023
%     \cite{iso25010}, y CUANTIFICADOS con una metrica.
%   * Cada requisito debe ser singular (un solo "shall" verificable),
%     no ambiguo y verificable. Sin terminos prohibidos por 5.2.7.
%   * La columna ISO/IEC 25010 solo se llena en los requisitos de calidad;
%     en los demas se deja con una raya.
%
% Referencia rapida del modelo de producto de ISO/IEC 25010:2023
% (verifica nombres y numeracion en el documento oficial antes de citar;
% en la version 2023 "usabilidad" paso a capacidad de interaccion y
% "portabilidad" paso a flexibilidad, y se agrego seguridad fisica):
%   - Idoneidad funcional: completitud, correccion, pertinencia.
%   - Eficiencia de desempeno: comportamiento temporal, utilizacion de
%     recursos, capacidad.
%   - Compatibilidad: coexistencia, interoperabilidad.
%   - Capacidad de interaccion: reconocibilidad, aprendibilidad,
%     operabilidad, proteccion contra errores, atractivo, inclusion,
%     asistencia, autodescripcion.
%   - Fiabilidad: ausencia de fallos, disponibilidad, tolerancia a
%     fallos, capacidad de recuperacion.
%   - Seguridad: confidencialidad, integridad, no repudio,
%     responsabilidad, autenticidad, resistencia.
%   - Mantenibilidad: modularidad, reutilizabilidad, analizabilidad,
%     modificabilidad, capacidad de prueba.
%   - Flexibilidad: adaptabilidad, escalabilidad, instalabilidad,
%     reemplazabilidad.
%   - Seguridad fisica: restriccion operacional, identificacion de
%     riesgos, a prueba de fallos, advertencia de peligro, integracion
%     segura.

Los quince requisitos se presentan en la Tabla~\ref{tab:requisitos}.

{\footnotesize
\setlength{\tabcolsep}{3pt}
\begin{longtable}{@{}p{1.0cm}p{4.6cm}p{1.9cm}p{2.2cm}p{2.0cm}p{2.2cm}@{}}
\caption{Requisitos del sistema, tipo, clasificación de calidad,
métrica y criterio de aceptación.}
\label{tab:requisitos}\\
\toprule
\textbf{ID} & \textbf{Enunciado} & \textbf{Tipo (5.2)} &
\textbf{ISO/IEC 25010:2023} & \textbf{Métrica} &
\textbf{Criterio de aceptación} \\
\midrule
\endfirsthead
\multicolumn{6}{@{}l}{\footnotesize\itshape Continuación de la Tabla~\ref{tab:requisitos}}\\
\toprule
\textbf{ID} & \textbf{Enunciado} & \textbf{Tipo (5.2)} &
\textbf{ISO/IEC 25010:2023} & \textbf{Métrica} &
\textbf{Criterio de aceptación} \\
\midrule
\endhead
\bottomrule
\endfoot
R01 & \pd & \pd & \pd & \pd & \pd \\
R02 & \pd & \pd & \pd & \pd & \pd \\
R03 & \pd & \pd & \pd & \pd & \pd \\
R04 & \pd & \pd & \pd & \pd & \pd \\
R05 & \pd & \pd & \pd & \pd & \pd \\
R06 & \pd & \pd & \pd & \pd & \pd \\
R07 & \pd & \pd & \pd & \pd & \pd \\
R08 & \pd & \pd & \pd & \pd & \pd \\
R09 & \pd & \pd & \pd & \pd & \pd \\
R10 & \pd & \pd & \pd & \pd & \pd \\
R11 & \pd & \pd & \pd & \pd & \pd \\
R12 & \pd & \pd & \pd & \pd & \pd \\
R13 & \pd & \pd & \pd & \pd & \pd \\
R14 & \pd & \pd & \pd & \pd & \pd \\
R15 & \pd & \pd & \pd & \pd & \pd \\
\end{longtable}
}

% =====================================================================
\section{Autoevaluación y corrección de los requisitos}
\label{sec:autoevaluacion}

% GUIA (borrar al terminar).
%   * Revisa CADA requisito contra las nueve caracteristicas de 5.2.5 y
%     contra los terminos prohibidos de 5.2.7.
%   * Para los que fallen, deja la version anterior y la corregida en la
%     tabla, y explica en una linea que caracteristica fallaba.
%   * Minimo recomendado: tres requisitos reescritos.
%
% Lista de control de las nueve caracteristicas de 5.2.5:
%   necesaria, apropiada, no ambigua, completa, singular, factible,
%   verificable, correcta, conforme.
%
% Terminos y expresiones a buscar en tus enunciados (5.2.7): "etc.",
% "y/o", "apropiado", "adecuado", "según corresponda", "entre otros",
% "soporta", "amigable", "eficiente", "flexible", "robusto", "rapido",
% "facil", "moderno", "óptimo".

Los resultados de la revisión se resumen en la
Tabla~\ref{tab:autoevaluacion}.

{\small
\setlength{\tabcolsep}{3pt}
\begin{longtable}{@{}p{1.0cm}p{2.9cm}p{5.25cm}p{5.25cm}@{}}
\caption{Requisitos que no superaron la revisión: versión anterior y
versión corregida.}
\label{tab:autoevaluacion}\\
\toprule
\textbf{ID} & \textbf{Característica o término que falla} &
\textbf{Versión anterior} & \textbf{Versión corregida} \\
\midrule
\endfirsthead
\multicolumn{4}{@{}l}{\footnotesize\itshape Continuación de la Tabla~\ref{tab:autoevaluacion}}\\
\toprule
\textbf{ID} & \textbf{Característica o término que falla} &
\textbf{Versión anterior} & \textbf{Versión corregida} \\
\midrule
\endhead
\bottomrule
\endfoot
\pd & \pd & \pd & \pd \\
\pd & \pd & \pd & \pd \\
\pd & \pd & \pd & \pd \\
\end{longtable}
}

% =====================================================================
\section{Reflexión}
\label{sec:reflexion}

% GUIA (borrar al terminar). Media página, en primera persona y con
% ejemplos concretos de TU trabajo:
%   * Qué tipo de requisito te costó más formular y por qué (no digas
%     "fue difícil": explica qué decisión no sabías tomar).
%   * Qué error cometiste al principio y cómo lo detectaste al aplicar
%     5.2.5 o 5.2.7.
%   * Qué cambiarías en el proceso si volvieras a empezar.

\pendiente{: reflexión de media página sobre la dificultad concreta al
formular los requisitos y su causa.}

% =====================================================================
\section{Declaración de uso de inteligencia artificial}
\label{sec:ia}

% GUIA (borrar al terminar). La guia exige declarar herramienta,
% version, proposito y forma de verificacion. Si no usaste IA, deja una
% declaracion expresa de no uso. No borres esta seccion.

Se declara el uso de \pendiente{: nombre de la herramienta y versión}
como apoyo para \pendiente{: propósito concreto, por ejemplo: estructura
del documento en LaTeX y corrección de errores de compilación}. El
contenido sustantivo del informe (análisis, glosario, requisitos,
autoevaluación y reflexión) fue elaborado por el estudiante. Lo obtenido
se verificó mediante \pendiente{: cómo se verificó, por ejemplo:
compilación del documento, contraste con la cláusula citada de la norma
y revisión requisito por requisito}.

% =====================================================================
\printbibliography[title={Referencias}]

\end{document}
